% !TeX root = main.tex
\section{Metodología}
\label{sec:metodologia}

\subsection{Imagen de entrada e inicialización}

Se emplea una implementación de segmentación interactiva basada en GrabCut
para imágenes bidimensionales en escala de grises. El procedimiento alterna
la estimación de modelos de intensidad para el fondo y el primer plano con
la actualización de etiquetas mediante cortes de grafos. La descripción
siguiente corresponde a la implementación del cuaderno del
proyecto~\cite{cuaderno}, basada en los principios de GrabCut~\cite{rother2004}.
OpenCV se utiliza para leer las imágenes y transformar sus canales; la
segmentación se calcula mediante las funciones del cuaderno y no mediante
una llamada a \texttt{cv2.grabCut}.

Sea $\Omega$ el conjunto de píxeles de la imagen y $x_i\in[0,255]$ la
intensidad del píxel $i$. La imagen PNG se carga en escala de grises y las
intensidades se convierten a precisión doble para los cálculos. No se aplica
una normalización adicional, filtrado ni cambio de tamaño en el procedimiento
del cuaderno. La unidad de procesamiento es un corte; cada ejecución
selecciona un caso mediante su identificador.

La interacción consiste en delimitar sobre la imagen un rectángulo $R$
que contenga la lesión y un margen de tejido circundante, conforme al
esquema de inicialización rectangular de GrabCut~\cite{rother2004}. Sus coordenadas
se expresan en píxeles enteros. Se inicializan las etiquetas como
\begin{equation}
 \alpha_i^{(0)}=[i\in R],
 \label{eq:inicializacion}
\end{equation}
donde los corchetes valen uno cuando se cumple la condición y cero en caso
contrario. Las etiquetas pertenecen a $\{0,1\}$. La etiqueta cero representa fondo y la etiqueta uno, primer
plano. Durante toda la optimización se impone $\alpha_i=0$ para
$i\notin R$. Las etiquetas interiores pueden cambiar. Por tanto, el
rectángulo define el dominio de búsqueda y no la segmentación final. La
inicialización requiere muestras suficientes dentro y fuera del rectángulo
para formar los cinco grupos de cada modelo.

\subsection{Modelos de intensidad}

Para cada clase $a\in\{0,1\}$ se modela la intensidad mediante una mezcla
de $K=5$ distribuciones gaussianas unidimensionales. Esta es una
adaptación escalar del modelo de apariencia de GrabCut~\cite{rother2004}:
\begin{equation}
 p(x_i\mid a;\theta)=\sum_{k=1}^{K}w_{ak}
 \mathcal{N}(x_i;\mu_{ak},\sigma_{ak}^{2}),
 \label{eq:mezcla}
\end{equation}
donde $\theta$ reúne los pesos $w_{ak}$, las medias $\mu_{ak}$ y las
varianzas $\sigma_{ak}^{2}$. Los pesos de cada clase satisfacen
\begin{equation}
 \sum_{k=1}^{K}w_{ak}=1.
 \label{eq:normalizacion-pesos}
\end{equation}
En la implementación del cuaderno~\cite{cuaderno}, la inicialización
ordena las intensidades de
cada clase y las divide en cinco grupos de tamaño aproximadamente igual.
El peso inicial de cada componente es la fracción de muestras de su grupo.
Su media y varianza se calculan a partir de esas muestras, añadiendo
$\varepsilon=10^{-6}$ a la varianza.

En cada iteración, cada píxel se asigna al componente de mayor contribución
ponderada dentro de su clase actual, siguiendo el esquema de asignación
única de componentes~\cite{rother2004}. De forma equivalente,
\begin{equation}
 k_i=\operatorname*{arg\,min}_{k}d_{\alpha_i k}(x_i),
\end{equation}
donde el costo de cada componente es
\begin{equation}
 d_{ak}(x_i)=-\log w_{ak}+\tfrac12\log(2\pi\sigma_{ak}^{2})
 +\frac{(x_i-\mu_{ak})^2}{2\sigma_{ak}^{2}}.
\end{equation}
El logaritmo es natural. Los componentes de peso cero tienen costo infinito.
La asignación es única y no modifica todavía la clase del píxel.

Sean $X_{ak}$ las muestras asignadas al componente $k$ de la clase $a$,
$N_{ak}=|X_{ak}|$ y $N_a=\sum_k N_{ak}$. Los parámetros se actualizan como
\begin{equation}
 w_{ak}=\frac{N_{ak}}{N_a},
 \label{eq:actualizacion-pesos}
\end{equation}
\begin{equation}
 \mu_{ak}=\frac{1}{N_{ak}}\sum_{x_i\in X_{ak}}x_i,
 \label{eq:actualizacion-medias}
\end{equation}
\begin{equation}
 \sigma_{ak}^{2}=\frac{1}{N_{ak}}\sum_{x_i\in X_{ak}}(x_i-\mu_{ak})^2
 +\varepsilon.
 \label{eq:actualizacion-modelos}
\end{equation}
Un componente sin muestras recibe peso cero y conserva su media y varianza
anteriores. Si una clase queda vacía, se conservan todos sus parámetros.
El costo de asignar el píxel a una clase se calcula con la mezcla
completa~\cite{cuaderno}:
\begin{equation}
 D_i(a;\theta)=-\log p(x_i\mid a;\theta).
 \label{eq:costo-apariencia}
\end{equation}
La suma se evalúa en el dominio logarítmico para reducir problemas de
precisión numérica. Esta elección difiere de la energía con componentes
explícitos del artículo original~\cite{rother2004}: la
asignación de componentes se utiliza para ajustar los parámetros, pero el
costo del corte marginaliza sobre todos los componentes.

\subsection{Energía y coherencia espacial}

La segmentación combina el costo de intensidad con una penalización por
asignar etiquetas diferentes a píxeles vecinos, siguiendo la formulación
de segmentación por grafos~\cite{boykov2001,rother2004}:
\begin{equation}
 E(\alpha;\theta)=\sum_{i\in\Omega}D_i(\alpha_i;\theta)
 +\sum_{\{i,j\}\in\mathcal{N}}v_{ij}[\alpha_i\ne\alpha_j].
 \label{eq:energia-segmentacion}
\end{equation}
El conjunto $\mathcal{N}$ contiene los pares de una vecindad de ocho
vecinos, contando cada par una sola vez. Se emplean los pesos espaciales
y la estimación de contraste descritos para imágenes escalares
en~\cite{rother2004}:
\begin{equation}
 v_{ij}=\frac{\gamma}{d_{ij}}\exp[-\beta(x_i-x_j)^2],
 \label{eq:pesos-espaciales}
\end{equation}
donde
\begin{equation}
 \beta=\left(\frac{2}{|\mathcal{N}|}
 \sum_{\{i,j\}\in\mathcal{N}}(x_i-x_j)^2\right)^{-1},
 \label{eq:contraste-espacial}
\end{equation}
con $\gamma=10$, valor fijado en el cuaderno~\cite{cuaderno}. Este valor
es una elección de la implementación y no el ajuste $\gamma=50$
reportado en el artículo original~\cite{rother2004}. La distancia $d_{ij}$ vale uno para vecinos horizontales
y verticales, y $\sqrt{2}$ para diagonales, en unidades de píxel. El
parámetro $\beta$ se estima sobre toda la imagen; si todas las diferencias
son nulas, se fija en cero. Los pesos elevados penalizan separar píxeles
de intensidad similar, mientras que las diferencias grandes reducen la
penalización. Estos pesos permanecen fijos durante las iteraciones.

\subsection{Construcción del grafo y corte mínimo}

La representación mediante terminales y enlaces entre vecinos sigue el
principio de los cortes de grafos para segmentación~\cite{boykov2001}.
En el cuaderno, el grafo incluye únicamente los píxeles de $R$ como nodos
variables~\cite{cuaderno}. El
efecto de los vecinos exteriores, fijados como fondo, se incorpora al costo
de primer plano mediante
\begin{equation}
 q_i=\sum_{\substack{j\notin R\\\{i,j\}\in\mathcal{N}}}v_{ij},
\end{equation}
\begin{equation}
 \widetilde D_i(a;\theta)=D_i(a;\theta)+a q_i.
\end{equation}
Así, la energía que depende de las etiquetas interiores es
\begin{equation}
 E_R(\alpha;\theta)=\sum_{i\in R}\widetilde D_i(\alpha_i;\theta)
 +\sum_{\substack{\{i,j\}\in\mathcal{N}\\i,j\in R}}
 v_{ij}[\alpha_i\ne\alpha_j].
 \label{eq:energia-reducida}
\end{equation}
Para modelos fijos, la diferencia entre $E$ y $E_R$ es el costo constante
$\sum_{i\notin R}D_i(0;\theta)$ y no cambia el corte óptimo.

Se añaden una fuente $s$, asociada al primer plano, y un sumidero $t$,
asociado al fondo. Para obtener capacidades no negativas, se resta de ambos
costos de cada píxel el valor
$m_i=\min\{\widetilde D_i(0;\theta),\widetilde D_i(1;\theta)\}$.
Las capacidades terminales son
\begin{equation}
 c(s,i)=\widetilde D_i(0;\theta)-m_i,
\end{equation}
\begin{equation}
 c(i,t)=\widetilde D_i(1;\theta)-m_i.
\end{equation}
Cada par de vecinos interiores se conecta mediante dos arcos opuestos de
capacidad $v_{ij}$. Si las etiquetas difieren, el costo del corte incluye
una sola contribución $v_{ij}$.

El flujo máximo se calcula mediante el algoritmo de Dinic~\cite{dinic1970},
con tolerancia
residual de $10^{-10}$ fijada en la implementación~\cite{cuaderno}. El procedimiento alterna la construcción de niveles
en el grafo residual con el envío de flujo, hasta que el sumidero deja de
ser alcanzable desde la fuente. Los píxeles interiores que permanecen
alcanzables reciben la etiqueta uno; los demás reciben cero. Con los
modelos fijos, la equivalencia entre flujo máximo y corte mínimo permite
minimizar la energía binaria~\cite{boykov2001,szeliski2011} de la
ecuación~\eqref{eq:energia-reducida}, salvo tolerancia numérica. Su valor se
recupera sumando $\sum_{i\in R}m_i$ al flujo máximo.

\subsection{Iteración y registro}

La ejecución completa parte de las etiquetas de la
ecuación~\eqref{eq:inicializacion} y de modelos recién inicializados.
La ejecución configurada en el cuaderno realiza cinco
iteraciones~\cite{cuaderno}. En cada una se asignan los componentes,
se actualizan los parámetros gaussianos, se recalculan los costos y se
obtiene un nuevo corte mínimo. El exterior del rectángulo y los pesos
espaciales permanecen fijos. Se conserva la máscara de cada iteración y se
registra el número de etiquetas modificadas. Las cinco iteraciones se
ejecutan incluso si las etiquetas dejan de cambiar; no se utiliza una
condición de parada basada en convergencia.

Para cada iteración $\ell$, se comparan las energías anterior y posterior
al corte con los mismos modelos actualizados $\theta^{(\ell)}$:
\begin{equation}
 E_R(\alpha^{(\ell)};\theta^{(\ell)})
 \leq E_R(\alpha^{(\ell-1)};\theta^{(\ell)}).
\end{equation}
Esta propiedad corresponde al paso de corte. La estimación de parámetros
emplea asignaciones únicas de componentes, mientras que el costo de
clasificación utiliza la mezcla completa. Por ello, el reajuste de las
gaussianas no garantiza una disminución de esta energía entre iteraciones.
La máscara de la quinta iteración constituye el resultado final.

\subsection{Evaluación por corte}

La segmentación final $\widehat\alpha$ se compara con la referencia $g_i$
definida en la ecuación~\eqref{eq:referencia-roja}. Para un dominio de
evaluación $A\subseteq\Omega$, la métrica implementada en el
cuaderno~\cite{cuaderno} es el error cuadrático medio:
\begin{equation}
 \operatorname{ECM}_A=\frac{1}{|A|}
 \sum_{i\in A}(\widehat\alpha_i-g_i)^2.
 \label{eq:ecm}
\end{equation}
Se calcula para la imagen completa ($A=\Omega$) y para el rectángulo
($A=R$). Al tratarse de etiquetas binarias, cada discrepancia contribuye
con uno; el error equivale a la proporción de píxeles mal clasificados.
Además, se cuentan los falsos positivos y falsos negativos sobre toda la
imagen y los píxeles de referencia que quedan fuera de $R$. Estos últimos
no pueden recuperarse porque el exterior permanece fijado como fondo.

Se incluye como comparación una predicción que asigna fondo a todos los
píxeles. Esta comparación permite interpretar un error global pequeño
cuando la lesión ocupa una fracción reducida de la imagen. La inspección
visual reúne la anotación, la máscara de referencia, la segmentación y un
mapa de errores dentro del rectángulo. Las métricas globales incluyen los
errores exteriores que no aparecen en esa visualización.

\subsection{Disponibilidad del código y aplicación web}
\label{sec:disponibilidad}

Todo el código del proyecto se encuentra en el repositorio
\textit{horus-grabcut}~\cite{repositorio}, disponible en
\url{https://github.com/DarthFranco2001/horus-grabcut}.
El cuaderno de procesamiento se ubica en \path{notebooks/}; el código de
la aplicación web y sus pruebas, en \path{web/}; y las imágenes y
anotaciones utilizadas, en \path{data/}. La referencia al cuaderno
identifica la copia local de la que proceden las salidas de este reporte.

La aplicación \href{https://darthfranco2001.github.io/horus-grabcut/}{\textit{Laboratorio GrabCut}}~\cite{web} está disponible en línea.
Permite seleccionar un caso, delimitar el rectángulo de interés, ejecutar
la segmentación y consultar la comparación con la referencia. El
procesamiento se realiza en el navegador. La aplicación ofrece un medio
interactivo para explorar el procedimiento; los resultados cuantitativos
presentados en este reporte corresponden a la ejecución del cuaderno
descrita en la sección~\ref{sec:resultados}.
